% !TeX root = ../supplementary.tex
\section{Qwen2.5-VL Prompt and Decoding Protocol}
\label{sec:supp:qwen}

The three locked conditions use Qwen2.5-VL-7B-Instruct~\cite{bai2025qwen25vl}
at revision \nolinkurl{cc594898137f460bfe9f0759e9844b3ce807cfb5}, bfloat16,
SDPA, no quantization, and a slow processor with 200,704--1,003,520 allowed
image pixels.  Prompt templates, the label vocabulary, processor policy, and
decoding policy are content-hashed before target inference.  Prompt
development and parse checks use HarMeme only; the prompt-lock records report
zero FHM label reads and zero FHM model forwards before locking.

\textsc{Direct} requests the structured JSON prediction without a rationale.
\textsc{Rationale} requests the same structured output plus a concise,
evidence-grounded justification; we do not expose or evaluate hidden
chain-of-thought.  \textsc{RAG} additionally supplies evidence retrieved from
the frozen HarMeme-train-only resource.  Its result remains pending and no
value is included here.

Decoding is deterministic (sampling disabled, one beam, temperature unset,
and at most 384 new tokens).  One deterministic repair request is permitted
only after JSON parsing, required-field, type, or canonical-vocabulary
failure; it instructs the model to repair the schema without adding claims.
The formal binary decision is candidate-likelihood argmax over
\texttt{non\_harmful} and \texttt{harmful}, normalized by a two-candidate
softmax.  Self-reported confidence is not used as a score.
